\documentclass[11pt,aspectratio=169]{ctexbeamer}

\usetheme{default}
\setbeamertemplate{navigation symbols}{}

\usepackage{amsmath}
\usepackage{booktabs}

\definecolor{DeepBlue}{RGB}{25,74,122}
\definecolor{SoftInk}{RGB}{53,59,66}
\definecolor{PaperBg}{RGB}{249,248,245}
\definecolor{PanelBlue}{RGB}{236,243,249}

\setbeamercolor{normal text}{fg=SoftInk,bg=PaperBg}
\setbeamercolor{structure}{fg=DeepBlue}
\setbeamercolor{frametitle}{fg=DeepBlue,bg=PaperBg}
\setbeamercolor{title}{fg=DeepBlue}
\setbeamercolor{block title}{fg=DeepBlue,bg=PanelBlue}
\setbeamercolor{block body}{fg=SoftInk,bg=white}

\newcommand{\Core}[1]{
  \begin{block}{核心结论}
    #1
  \end{block}
}

\newcommand{\LectureMeta}[2]{
  \title{大模型推理基础设施}
  \subtitle{#1}
  \author{张书豪}
  \institute{华中科技大学计算机学院}
  \date{#2}
}


\LectureMeta{第 2 讲 工作负载与评价指标}{研究生课程课件}

\begin{document}

\begin{frame}[plain]
  \titlepage
\end{frame}

\begin{frame}{教学目标}
  \Core{本讲围绕工作负载与评价指标展开，建立刻画推理系统问题的基本分析坐标。}
  \begin{itemize}
    \item 区分离线、在线、RAG 与智能体服务。
    \item 理解 TTFT、TPOT、P99、goodput 等核心指标。
    \item 说明为什么单点指标无法代表端到端系统收益。
  \end{itemize}
\end{frame}

\begin{frame}{典型工作负载类型}
  \Core{不同工作负载会重新定义系统瓶颈，因此系统设计不能脱离场景。}
  \begin{itemize}
    \item 离线批推理：关注总吞吐与资源利用率。
    \item 在线对话：关注 TTFT、TPOT 和尾时延。
    \item RAG 服务：关注检索到生成的全链路时延。
    \item 智能体工作流：关注跨轮状态、工具调用与任务级完成时间。
  \end{itemize}
\end{frame}

\begin{frame}{四类工作负载的系统差异}
  \Core{同样叫“推理”，不同 workload 的主导矛盾差别很大。}
  \begin{itemize}
    \item 离线推理更偏向吞吐和成本优化。
    \item 在线对话更偏向 TTFT 和尾时延。
    \item RAG 更偏向检索与生成链路的协同。
    \item 智能体更偏向多阶段流水和跨轮状态维护。
  \end{itemize}
\end{frame}

\begin{frame}{为什么工作负载变了，系统问题也会变}
  \Core{长上下文、MoE、多模态和智能体会从容量、调度、通信和状态维护多个层面放大复杂度。}
  \begin{itemize}
    \item 长上下文放大 prefill 和 KV 压力。
    \item MoE 放大负载倾斜与通信波动。
    \item 智能体把单轮生成变成多阶段流水与持续写入问题。
  \end{itemize}
\end{frame}

\begin{frame}{长上下文如何改变瓶颈}
  \Core{长上下文会同时推高 prefill 计算代价、KV 驻留压力与排队放大效应。}
  \begin{itemize}
    \item prompt 变长直接增加 prefill 开销。
    \item KV 驻留时间更长，占用更大。
    \item 长短请求混部时更容易出现尾时延恶化。
  \end{itemize}
\end{frame}

\begin{frame}{核心评价指标}
  \Core{评价指标的选择实质上对应系统优化目标的定义。}
  \begin{itemize}
    \item TTFT：首 token 等待时间。
    \item TPOT：每个输出 token 的平均间隔。
    \item Throughput：单位时间完成的请求或 token 数。
    \item P95/P99：尾时延表现。
    \item Goodput：满足服务目标的有效吞吐。
  \end{itemize}
\end{frame}

\begin{frame}{指标之间为什么会冲突}
  \Core{优化一个指标常常会伤害另一个指标，因此系统研究首先要明确主目标和约束。}
  \begin{itemize}
    \item 更大的 batch 通常提高吞吐，但可能恶化 TTFT。
    \item 更激进缓存策略可能提高命中率，但也提高管理代价。
    \item 更高设备利用率不等于更低 P99。
  \end{itemize}
\end{frame}

\begin{frame}{为什么平均值常常不够}
  \Core{用户体验和服务质量常由尾部与波动决定，而不是平均值。}
  \begin{itemize}
    \item 更大的 batch 可能提高平均吞吐，但也可能恶化 TTFT。
    \item 更高命中率未必自动带来端到端时延下降。
    \item 单次最好结果并不等于稳定收益。
  \end{itemize}
\end{frame}

\begin{frame}{为什么要引入 goodput}
  \Core{goodput 的意义在于把“服务目标是否满足”纳入吞吐定义，而不是只统计原始完成量。}
  \begin{itemize}
    \item 如果大量请求违反 SLO，原始 throughput 的意义会迅速下降。
    \item goodput 更适合比较不同调度或资源治理策略。
    \item 它把系统目标从“跑得多”转化为“有效跑得多”。
  \end{itemize}
\end{frame}

\begin{frame}{讲解：工作负载为什么先于指标}
  \Core{指标没有脱离场景的独立意义，必须先交代 workload，才能解释数字。}
  \begin{itemize}
    \item 离线批处理更接近资源利用率问题，在线对话更接近响应性问题。
    \item RAG 与 Agent 会把“生成”嵌入更长链路，因此单点吞吐不再足以代表体验。
    \item 同一套指标在不同 workload 下含义会变化，这就是为什么课程里反复强调“先说场景，再说指标”。
  \end{itemize}
\end{frame}

\begin{frame}{讲解：这一讲真正要学会的不是指标名词}
  \Core{这讲的核心不是记住 TTFT、P99、goodput 的定义，而是学会用它们刻画取舍。}
  \begin{itemize}
    \item 一个系统如果吞吐更高，但首 token 更慢，它未必更适合在线服务。
    \item 一个系统如果尾时延稳定，但设备空转严重，它也未必更优。
    \item 因此评价不是找“最好看的数字”，而是明确哪个指标是主目标，哪个指标是代价。
  \end{itemize}
\end{frame}

\begin{frame}{例子：同样是在线请求，工作负载分布可以完全不同（1）}
  \Core{指标解释必须建立在真实 workload 之上，否则结论会整体偏移。}
\begin{block}{对话服务}
    请求到达更连续，prompt 长度波动大，用户对 TTFT 更敏感。
  \end{block}
\end{frame}

\begin{frame}{例子：同样是在线请求，工作负载分布可以完全不同（2）}
  \Core{指标解释必须建立在真实 workload 之上，否则结论会整体偏移。}
\begin{block}{RAG 服务}
    请求链路更长，检索与生成交织，单次输出长度未必长，但整体时延更复杂。
  \end{block}
\end{frame}

\begin{frame}{例子：同样是在线请求，工作负载分布可以完全不同（3）}
  \Core{指标解释必须建立在真实 workload 之上，否则结论会整体偏移。}
\begin{block}{Agent 服务}
    一次用户任务会触发多轮工具调用和多次生成，简单平均吞吐很难刻画真实体验。
  \end{block}
\end{frame}

\begin{frame}{例子：为什么平均值经常误导课程项目（1）}
  \Core{系统课里最常见的误区，是只看均值而不看分布和尾部。}
\begin{block}{平均 TTFT 看起来正常}
    但只要少量长请求和突发到达叠加，P99 就可能已经不可接受。
  \end{block}
\end{frame}

\begin{frame}{例子：为什么平均值经常误导课程项目（2）}
  \Core{系统课里最常见的误区，是只看均值而不看分布和尾部。}
\begin{block}{平均吞吐更高}
    可能只是牺牲了短请求体验或拉高了队列等待，并不等于服务更好。
  \end{block}
\end{frame}

\begin{frame}{例子：为什么平均值经常误导课程项目（3）}
  \Core{系统课里最常见的误区，是只看均值而不看分布和尾部。}
\begin{block}{goodput 的作用}
    它提醒我们：真正有价值的是满足 SLO 的有效吞吐，而不是单纯多算几个 token。
  \end{block}
\end{frame}

\begin{frame}{代码例子：最小 benchmark 如何固定输入分布（1）}
  \Core{工作负载与指标课程需要让学生看到，实验不是“随便跑一遍”。}
\begin{block}{代码片段}
    {\ttfamily\small
    \detokenize{seed(0)}\par
    \detokenize{prompt_token_ids = [[randint(0, 10000) for _ in range(randint(100, max_input_len))] ...]}\par
    \detokenize{sampling_params = [SamplingParams(max_tokens=randint(100, max_output_len)) ...]}\par
    }
  \end{block}
\end{frame}

\begin{frame}{代码例子：最小 benchmark 如何固定输入分布（2）}
  \Core{工作负载与指标课程需要让学生看到，实验不是“随便跑一遍”。}
\begin{block}{课堂应追问}
    这里哪些变量控制了“输入长度分布”，哪些变量控制了“输出长度分布”？
  \end{block}
\end{frame}

\begin{frame}{代码例子：最小 benchmark 如何固定输入分布（3）}
  \Core{工作负载与指标课程需要让学生看到，实验不是“随便跑一遍”。}
\begin{block}{系统含义}
    如果这两类分布没说清，后面的 TTFT、throughput、queue wait 都可能失去解释基础。
  \end{block}
\end{frame}

\begin{frame}{代码例子：如何把原始运行结果转成指标（1）}
  \Core{指标不是抽象名词，而是从运行记录里计算出来的。}
\begin{block}{最小计算逻辑}
    {\ttfamily\small
    \detokenize{t = time.time()}\par
    \detokenize{llm.generate(prompt_token_ids, sampling_params, use_tqdm=False)}\par
    \detokenize{throughput = total_tokens / (time.time() - t)}\par
    }
  \end{block}
\end{frame}

\begin{frame}{代码例子：如何把原始运行结果转成指标（2）}
  \Core{指标不是抽象名词，而是从运行记录里计算出来的。}
\begin{block}{局限}
    这里只有总体 throughput，没有拆 TTFT、P99 和批次级探针。
  \end{block}
\end{frame}

\begin{frame}{代码例子：如何把原始运行结果转成指标（3）}
  \Core{指标不是抽象名词，而是从运行记录里计算出来的。}
\begin{block}{教学意义}
    这正好说明：一个最小 benchmark 足够入门，但不足以支撑完整系统结论。
  \end{block}
\end{frame}

\begin{frame}{例子：throughput、TTFT、P99 三者为何经常互相拉扯（1）}
  \Core{工作负载与评价指标这一讲的重点，是学会接受真实系统中的指标张力。}
\begin{block}{更大的 batch}
    往往提高总体吞吐，但也可能拉高等待时间和 TTFT。
  \end{block}
\end{frame}

\begin{frame}{例子：throughput、TTFT、P99 三者为何经常互相拉扯（2）}
  \Core{工作负载与评价指标这一讲的重点，是学会接受真实系统中的指标张力。}
\begin{block}{更保守的调度}
    可能稳定尾时延，却压低设备利用率和平均吞吐。
  \end{block}
\end{frame}

\begin{frame}{例子：throughput、TTFT、P99 三者为何经常互相拉扯（3）}
  \Core{工作负载与评价指标这一讲的重点，是学会接受真实系统中的指标张力。}
\begin{block}{课程项目要求}
    任何优化都要先说清：你优先保哪个指标，愿意牺牲哪个指标。
  \end{block}
\end{frame}

\begin{frame}{课堂练习：为同一系统写出两种不同的指标故事（1）}
  \Core{同一组数字可以对应完全不同的解释，关键在于你把问题定义成什么。}
\begin{block}{练习一}
    如果一个系统吞吐提高 15\%，但 TTFT 恶化 30\%，你会如何判断它对在线服务是否更好？
  \end{block}
\end{frame}

\begin{frame}{课堂练习：为同一系统写出两种不同的指标故事（2）}
  \Core{同一组数字可以对应完全不同的解释，关键在于你把问题定义成什么。}
\begin{block}{练习二}
    给出一个适用于 Agent 工作流的 SLO 组合，不要只写单个 TTFT 阈值。
  \end{block}
\end{frame}

\begin{frame}{课堂练习：为同一系统写出两种不同的指标故事（3）}
  \Core{同一组数字可以对应完全不同的解释，关键在于你把问题定义成什么。}
\begin{block}{练习三}
    说明为什么课程里反复强调“指标必须回到 workload 和场景”。
  \end{block}
\end{frame}


\begin{frame}{课堂讨论}
  \Core{指标选择本身构成系统问题定义的重要组成部分。}
  \begin{itemize}
    \item 为什么以 TTFT 为核心的服务不会只优化峰值吞吐？
    \item RAG 服务和对话服务的主指标为何不同？
    \item 什么情况下应该报告 goodput 而不是 raw throughput？
  \end{itemize}
\end{frame}

\begin{frame}{课堂练习}
  \Core{本讲可组织“场景-指标-瓶颈”三列表练习，以训练问题刻画能力。}
  \begin{itemize}
    \item 每个场景至少写出 2 个主指标。
    \item 每个指标至少写出 1 个可能冲突指标。
    \item 每个冲突都要对应一个系统机制取舍。
  \end{itemize}
\end{frame}

\begin{frame}{本讲总结}
  \Core{明确工作负载与评价指标，是后续讨论调度、KV 与执行优化问题的共同前提。}
  \begin{itemize}
    \item 不同场景对应不同瓶颈。
    \item 指标定义决定优化方向和论文叙事边界。
  \end{itemize}
\end{frame}

\end{document}